\section{Pagan Philosophers and Church Fathers}

\begin{quotationx}
The precise nature and even, in some senses, the width of the chasm which separated the Pagan and Christian religions can easily be mistaken if we take our ideas solely from political or ecclesiastical histories: still more, if we take them from more popular sources. Cultured people on both sides had had the same education, read the same poets, learned the same rhetoric.

I have read a novel which represents all the Pagans of that day as carefree sensualists, and all the Christians as savage ascetics. It is a grave error. They were in some ways far more like each other than either was like a modern man. The leaders on both sides were monotheists, and both admitted almost an infinity of supernatural beings between God and man. Both were highly intellectual ... The last champions of Paganism were not the sort of men that Swinburne, or a modern ``humanist', would wish them to have been. They were not lusty extroverts recoiling in horror or contempt from a world grown grey with the breath of the `pale Galilean'. If they wanted to get back `the laurel, the palms, and the paean', it was on most serious and religious grounds. If they longed to see `the breast of the nymph in the brake', their longing was not like a satyr's; it was much more like a spiritualist's. A world-renouncing, ascetic, and mystical character then marked the most eminent Pagans no less than their Christian opponents. It was the spirit of the age. Everywhere, on both sides, men were turning away from the civic virtues and the sensual pleasures to seek an inner purgation and a supernatural goal. The modern who dislikes the Christian Fathers would have disliked the Pagan philosophers equally, and for similar reasons. Both alike would have embarrassed him with stories of visions, ecstasies, and apparitions. Between the lower and more violent manifestations of both religions he would have found it hard to choose. To a modern eye (and nostril) Julian with his long nails and densely populated beard might have seemed very like an unwashed monk out of the Egyptian desert.

\flright{\textsc{C. S. Lewis}, \textit{The Discarded Image}}
\end{quotationx}






\flrightit{Posted on 2009-12-23 by Cologero}

\begin{center}* * *\end{center}

\begin{footnotesize}
\begin{sffamily}
\texttt{GFJF on 2010-01-07 at 05:41 said:}

Rather than monotheists, I think it would be better to call the pagan philosophers monists. But what period precisely is he referring to? 400-600 AD? Plato's Academy and Byzantium?

\hfill

\texttt{Cologero on 2010-02-05 at 10:21 said:}

The quote is from a discussion around 400-600 AD, as paganism gave way to Christianity. The point is that they were closer to monotheism, rather than polytheism, as the neo-pagans of today would like to believe. I don't see how ``monist'' is a better description.

\hfill

\texttt{GFJF on 2010-02-12 at 04:04 said:}

``The point is that they were closer to monotheism, rather than polytheism, as the neo-pagans of today would like to believe.''

Okay. Allow me to be pedantic and explain what I meant. The pagans of that period were neoplatonists, by and large, and so it is right to call them monists, believing as they did in a supreme unity. They put the Demiurge in with the ``infinity of beings above man '', not in the ultimate position. You see where I'm coming from? But let's not pursue it. I'm fine with calling them monotheists.

\hfill

\texttt{Cologero on 2010-02-12 at 22:38 said:}

Gabeesh. And I'm fine with calling them neoplatonists since for many readers, ``monotheism'' has the connotation of an ``old man in the sky''. And Lewis' point was that the Christians --- at least the educated and philosophic ones --- were close in spirit and lifestyle to those pagans. I understand, too, that despite the similarities, there was an intellectual battle between the two groups for the hearts of minds of the people. The pagans eventually lost, in no small measure because their ideas were absorbed into the Christian system. Unfortunately --- at least from my point of view --- the Reformation succeeded in its attempt to re-Judaize Christianity and purge it of its Greek and Roman elements.

\hfill

\texttt{GFJF on 2010-02-16 at 10:41 said:}

Speaking of Christianity absorbing pagan ideas, I have the works of Pseudo Dionysius on order... which I'm eager to delve into. In that connection, would it be fair to group the latterday pagans in, spiritually, with the the sacerdotal, as opposed to the royal tradition? It seems to me that that would be one possible conclusion to draw from Lewis here.

\hfill

\texttt{Cologero on 2010-02-20 at 19:20 said:}

Latterday pagans? There is no latterday pagan Tradition, but I plan to address the issue of the religious sense appropriate to different castes in the coming week.

As for Dionysus, that is too complex to address in a comment. As you read him, keep in mind one thing that may make the experience richer: According to Rene Guenon, the angelic heirarchies represent higher, or supr-individual, states of the Being

\hfill

\texttt{Gornahoor on 2019-12-23 at 14:30 said:}

Most neopagans look back on the Middle Ages and find the dominant religion to be alien; they then ascribe it to sources in the Middle East. Actually, as Rene Guenon pointed out, the entire Middle Ages are alien to the modern mind. It takes an effort to put oneself into that mindset. So even the pagans are misunderstood by neopagans today.

\hfill

\texttt{Michael on 2019-12-24 at 09:57 said:}

To that point it seems strange that Neopagans promote the respect of ancestors as a value that they wish to pursue, yet they disrespect the choices of said ancestors with conversion (accept vs. reject). Of course many point out ``forced'' conversions, which doesn't make much sense at all as a real objection, many neopagans don't consider the reasons or conditions to why conversions were ``forced'' in certain areas and not in others or in certain individuals and not in others. That point doesn't even bring into the account that all lived and died based in their ``Pagan values'' regardless, all things will come into life and fade into death, rebirth seemed more confusing to the Pagans and many saw the answer was somewhere else.

Do we or the neopagans really stop to think and see how the Pagan philosophers when approached by something not yet understood by them, willingly accepted and saw for themselves? or when they fought against it?

\hfill

\texttt{Max on 2019-12-31 at 19:25 said:}

Lewis claims that the difference between Pagans and Christians are most marked in ``political or ecclesiastical histories'', and this is entirely true, since the difference concerns external modes of expression. Christianity came to represent a new form of politial organisation. The rise of philosophy itself was a transitional phenomena, paralleling developments among the Hebrews. These changes were in the air regardless of religion, or as Lewis says: ``everywhere, on both sides, men were turning away from the civic virtues''. For example, Stoics and Christians alike were cosmopolitans. On the other hand, the ecumene was the empire. The challenge amounts to understanding the origin and aim of these political reforms, how the changes took place, what it all means. While things does not happen by chance, causes are often obscure. Reading political and ecclesiastical history provides a feel for these changes, and what it means going forward, in terms of options available to us. The past demarcates the future. We need to understand our trajectory, in addition to the prevailing winds, in order to to maneouver the ship. Those are givens, which none of us can change regardless of personal opinion.

I was gifted this tie with a few words of wisdom set over a nautical chart:

``A mariner does not ask for fair wind. He learns to sail.''

For those who belong to the earth, just stay put, and we will bring back the news from across the ocean.


\hfill

\end{sffamily}
\end{footnotesize}

